\section{Discussion}
\label{sec:discussion}

GradTDDFT is designed to make the response kernel part of the
trainable computational object rather than an after-the-fact
diagnostic. This distinction is important for neural
exchange-correlation functionals because excitation energies depend
not only on the converged KS orbitals and eigenvalue gaps, but also
on the derivative structure of the learned energy density. The
energy-consistent local response construction therefore enforces
consistency among the neural XC energy, the XC potential, and the
local $f_{\rm xc}$ kernel.
This consistency is a central requirement when the same functional is
used in both SCF optimization and TDDFT response.

The evaluation first separates numerical correctness from model
capability. Matched PySCF calculations test whether the
differentiable implementation reproduces conventional TDA and full
TDDFT in the non-neural limit. After this check, the model-capability
benchmarks separate sources of error that are often mixed in
neural-functional studies. H$_2$, H$_2^+$, and N$_2$ ground-state
dissociation curves test whether the learned functional remains
stable inside an SCF loop across equilibrium and stretched
geometries. H$_2$ and H$_2^+$ first excited-state dissociation curves
then test whether the same functional gives compatible ground-state
surfaces and response kernels. The 50-molecule QM9
ground- and excited-state training set tests whether the
differentiable SCF--TDDFT training graph transfers from controlled
bond scans to chemically diverse small organic molecules under
separate total-energy and vertical-excitation objectives. This
separation is necessary because a low energy error, a stable SCF
solution, and an accurate excitation spectrum are related but not
equivalent requirements.

This is also why the present training protocol keeps ground-state and
excited-state supervision separate. A joint loss is mathematically
straightforward in a differentiable SCF--TDDFT graph, but it is not a
robust practical objective for the current neural functional. The
ground-state loss primarily drives the self-consistent density,
orbital energies, and total-energy surface, whereas the excited-state
loss drives the response kernel, excitation energies, oscillator
strengths, and state ordering. When both are optimized at the same
time, their gradients can compete: updates that improve the
ground-state surface may degrade the response kernel, and updates that
fit a target excitation can destabilize the SCF fixed point or damage
the ground-state energy landscape. In our training design, separating
the objectives makes these failure modes observable rather than hiding
them inside an unstable multi-target optimization.

The nonlocal channels define a separate theoretical boundary. For
semilocal channels, automatic differentiation of the local energy
density gives a direct route to a consistent response kernel. For the
HF channel, the current excited-state path uses the learned local
exchange field to define a scalar hybrid fraction and then applies the
ordinary exact-exchange response matrix. For the MP2/PT2 channel, the
response problem remains a singles TDA/Casida problem and the learned
correlation coefficient scales a root-specific CIS(D)-type doubles
correction. This separation is useful for charge-transfer states,
double-excitation character, and open-shell systems where the chosen
excited-state correction can change the qualitative state ordering.

Several extensions remain natural next steps. A future direct
orbital-Hessian treatment of the MP2/PT2 channel would provide a more
direct connection between dynamical correlation descriptors and
excited-state response. A
JAX-native integral path would reduce the separation between the fixed
molecular data layer and the differentiable SCF layer, although
efficient production use still requires optimized J/K construction
and memory management. General unrestricted differentiable SCF would
broaden the training targets to radicals, charged species, and
spin-polarized response calculations. Larger basis sets
and larger molecules will further test whether the response-kernel
construction remains practical when the orbital and quadrature spaces
grow.
